\section{Cambiar de modalidad: borrar, dejar en pausa y retomar}
Solo puede haber un trámite activo. Al final de tu trámite está \emph{¿Elegiste otra modalidad o quieres dejar este trámite?}. Si nadie más ha intervenido, el trámite \textbf{se borra}; si ya intervino alguien (tu director firmó, Apoyo a Grado validó…), \textbf{queda en pausa} en \emph{Abandonadas}, conserva todo y lo retomas con \boton{Retomar}.
\captura{comun_egresado_01_dejar_pliegue}{Al final de tu trámite: «¿Elegiste otra modalidad o quieres dejar este trámite?».}
\captura{comun_egresado_02_borrar}{Si nadie más intervino, el trámite se borra definitivamente; la pantalla lo explica y pide confirmarlo.}
\captura{comun_egresado_03_una_activa}{Solo un trámite activo: antes de comenzar otro, la plataforma explica qué pasará con el que tienes (aquí, quedará en pausa porque tu director ya firmó).}
\captura{comun_egresado_04_abandonadas}{«Abandonadas»: el trámite en pausa conserva todo y se retoma desde aquí.}
\captura{comun_egresado_05_retomar}{Retomar: confirmas; el trámite activo se borra o queda en pausa, según el caso.}
\captura{comun_egresado_06_pausa}{Con intervención de otras personas, el trámite no se borra: queda en pausa y se explica por qué.}
